\subsection{Lý thuyết cần nhớ}

\subsubsection{Mệnh đề, mệnh đề chứa biến}

\begin{kienthuc}
\begin{itemize}
	\item \textbf{Mệnh đề:} Mệnh đề là một câu khẳng định đúng hoặc sai.
	
	\item Câu khẳng định đúng gọi là \textbf{mệnh đề đúng}.
	
	\item Câu khẳng định sai gọi là \textbf{mệnh đề sai}.
	
	\item Một mệnh đề không thể vừa đúng, vừa sai.
	
	\item Những mệnh đề liên quan đến toán học được gọi là \textbf{mệnh đề toán học}.
	
	\item \textbf{Mệnh đề chứa biến:} Là câu khẳng định chứa biến nhận giá trị trong một tập xác định nào đó và với mỗi giá trị cụ thể của biến thuộc tập xác định, ta được một mệnh đề.
\end{itemize}
\end{kienthuc}

\subsubsection{Mệnh đề phủ định}

\begin{kienthuc}
\begin{itemize}
	\item Cho mệnh đề $P$. Mệnh đề ``không phải $P$'' gọi là \textbf{mệnh đề phủ định của $P$}.
	
	\item Mệnh đề phủ định của $P$ kí hiệu là $\overline{P}$.
	
	\item Nếu $P$ đúng thì $\overline{P}$ sai; nếu $P$ sai thì $\overline{P}$ đúng.
\end{itemize}
\end{kienthuc}

\subsubsection{Mệnh đề kéo theo và mệnh đề đảo}

\begin{kienthuc}
\begin{itemize}
	\item Cho hai mệnh đề $P$ và $Q$.
	
	\item Mệnh đề ``Nếu $P$ thì $Q$'' gọi là \textbf{mệnh đề kéo theo}, kí hiệu:
	\[
	P\Rightarrow Q.
	\]
	
	\item Mệnh đề $P\Rightarrow Q$ còn được phát biểu là ``$P$ kéo theo $Q$'' hoặc ``Từ $P$ suy ra $Q$''.
	
	\item Mệnh đề $P\Rightarrow Q$ chỉ sai khi $P$ đúng và $Q$ sai.
	
	\item Trong mệnh đề $P\Rightarrow Q$:
	\begin{itemize}
		\item $P$ là \textbf{giả thiết}.
		\item $Q$ là \textbf{kết luận}.
	\end{itemize}
	
	\item \textbf{Mệnh đề đảo:} Cho mệnh đề $P\Rightarrow Q$. Khi đó, $Q\Rightarrow P$ gọi là \textbf{mệnh đề đảo} của $P\Rightarrow Q$.
\end{itemize}
\end{kienthuc}

\begin{emcothe}
\begin{itemize}
	\item Với mệnh đề $P\Rightarrow Q$, có thể phát biểu:
	\begin{itemize}
		\item $P$ là điều kiện đủ để có $Q$.
		\item $Q$ là điều kiện cần để có $P$.
	\end{itemize}
\end{itemize}
\end{emcothe}

\subsubsection{Mệnh đề tương đương}

\begin{kienthuc}
\begin{itemize}
	\item Cho hai mệnh đề $P$ và $Q$. Mệnh đề ``$P$ nếu và chỉ nếu $Q$'' gọi là \textbf{mệnh đề tương đương}.
	
	\item Kí hiệu:
	\[
	P\Leftrightarrow Q.
	\]
	
	\item Mệnh đề $P\Leftrightarrow Q$ đúng khi cả hai mệnh đề $P\Rightarrow Q$ và $Q\Rightarrow P$ cùng đúng.
	
	\item Khi đó:
	\[
	P\Leftrightarrow Q
	\]
	có thể đọc là ``$P$ khi và chỉ khi $Q$'', ``$P$ tương đương $Q$''.
\end{itemize}
\end{kienthuc}

\begin{emcothe}
\begin{itemize}
	\item Với $P\Leftrightarrow Q$, ta có thể phát biểu:
	\begin{itemize}
		\item $P$ là điều kiện cần và đủ để có $Q$.
		\item $P$ khi và chỉ khi $Q$.
		\item $P$ tương đương $Q$.
	\end{itemize}
\end{itemize}
\end{emcothe}

\subsubsection{Mệnh đề có kí hiệu $\forall$, $\exists$}

\begin{kienthuc}
\begin{itemize}
	\item \textbf{Mệnh đề chứa kí hiệu với mọi:}
	\[
	\forall x\in X,\ P(x).
	\]
	
	Đọc là: ``Với mọi $x$ thuộc $X$, $P(x)$''.
	
	\item Mệnh đề này đúng khi \textbf{tất cả} các giá trị của $x\in X$ đều làm cho phát biểu $P(x)$ đúng.
	
	\item Nếu tìm được \textbf{ít nhất một} giá trị $x\in X$ làm cho $P(x)$ sai thì mệnh đề trên sai.
	
	\item \textbf{Mệnh đề chứa kí hiệu tồn tại:}
	\[
	\exists x\in X,\ P(x).
	\]
	
	Đọc là: ``Tồn tại $x$ thuộc $X$ sao cho $P(x)$''.
	
	\item Mệnh đề này đúng khi tìm được \textbf{ít nhất một} giá trị $x\in X$ làm cho $P(x)$ đúng.
	
	\item Nếu tất cả các giá trị $x\in X$ đều làm cho $P(x)$ sai thì mệnh đề trên sai.
\end{itemize}
\end{kienthuc}

\subsubsection{Phủ định mệnh đề chứa $\forall$, $\exists$}

\begin{emcothe}
\begin{itemize}
	\item Phủ định của mệnh đề ``Với mọi $x\in X$, $P(x)$'' là:
	\[
	\overline{\left(\forall x\in X,\ P(x)\right)}
	\Leftrightarrow
	\exists x\in X,\ \overline{P(x)}.
	\]
	
	\item Phủ định của mệnh đề ``Tồn tại $x\in X$ sao cho $P(x)$'' là:
	\[
	\overline{\left(\exists x\in X,\ P(x)\right)}
	\Leftrightarrow
	\forall x\in X,\ \overline{P(x)}.
	\]
	
	\item Cần nhớ:
	\[
	\boxed{\overline{\forall}=\exists}
	\qquad\text{và}\qquad
	\boxed{\overline{\exists}=\forall}.
	\]
\end{itemize}
\end{emcothe}